\section{Contextualización y motivación}

\subsection{Problema}

La imagen médica es una fuente central de información diagnóstica, pero un
algoritmo convencional necesita ejemplos representativos y correctamente
etiquetados de cada entidad que deberá reconocer. En radiología, obtener esas
etiquetas exige tiempo experto, los hallazgos pueden ser sutiles, existe
variabilidad entre observadores y la distribución de patologías es muy
desigual. Además, un clasificador cerrado aprende las categorías disponibles y
puede responder con excesiva confianza ante un patrón no representado. Estas
condiciones motivan la detección de anomalías: modelar la normalidad y señalar
desviaciones sin exigir ejemplos de todas las enfermedades posibles
\cite{pang2021review,ruff2021unifying}.

La formulación de una clase es especialmente natural cuando abundan muestras
sanas y las anomalías son heterogéneas. Sea $x$ una radiografía y $f_\theta$ un
modelo ajustado con el conjunto normal $\mathcal{D}_{N}$. Se busca una función
$s(x)$ tal que
\[
  s(x_N) < s(x_A), \qquad x_N\sim p_N,\quad x_A\not\sim p_N,
\]
sin utilizar $x_A$ para optimizar $\theta$. El resultado no es un diagnóstico,
sino una puntuación de incompatibilidad con la distribución aprendida.

\subsection{Por qué reconstrucción generativa}

Un autoencoder comprime la entrada en una representación latente y aprende a
reconstruirla \cite{hinton2006reducing}. Si solo observa imágenes normales, se
espera que reconstruya peor una alteración y que el residuo sirva como señal de
anomalía. La hipótesis es atractiva, pero no está garantizada: una red con mucha
capacidad puede copiar también una entrada anómala; el error puede estar
dominado por contraste, bordes o adquisición, y una reconstrucción visualmente
nítida no implica mejor discriminación \cite{baur2021autoencoders}.

El VAE introduce una representación probabilística regularizada mediante la
divergencia de Kullback--Leibler \cite{kingma2014vae}. Esto puede producir un
espacio latente más continuo, aunque también reconstrucciones suavizadas. Las
GAN enfrentan generador y discriminador \cite{goodfellow2014gan}; aplicaciones
como AnoGAN y GANomaly trasladan ese aprendizaje adversarial a detección de
anomalías \cite{schlegl2017anogan,akcay2019ganomaly}. La literatura médica
incluye autoencoders adversariales, VAE en cascada y autoencoders enmascarados
\cite{zhang2022improved,guo2021cvad,georgescu2023masked}. No obstante, las
diferencias de datasets, preprocesado y métricas dificultan saber qué parte de
la mejora procede realmente del modelo.

\subsection{Pregunta de investigación y alcance}

La pregunta principal es: \emph{¿qué familia, AE, VAE o GANomaly, separa mejor
radiografías normales y no normales cuando todas se entrenan solo con normales
y se evalúan bajo idénticas condiciones?}

El estudio usa radiografías frontales del benchmark Chest-RSNA de BMAD
\cite{bao2024bmad,shih2019rsna}. ``Anómala'' significa \emph{no normal según la
taxonomía RSNA/BMAD}; no significa cualquier patología torácica ni confirma
neumonía. La salida podría servir como señal de priorización o control de
calidad en investigación, pero el sistema no está calibrado prospectivamente,
no ha sido validado externamente y no debe utilizarse para decisiones clínicas.

\subsection{Aportaciones}

Las aportaciones concretas son: (i) selección y auditoría reproducible de un
conjunto público adecuado; (ii) implementaciones comparables de AE, VAE y
GANomaly; (iii) separación estricta entre ajuste, selección de umbral y prueba;
(iv) evaluación con métricas robustas al desequilibrio y análisis de coste; y
(v) publicación del código y de la configuración, sin redistribuir imágenes
médicas.
